% =====================================================================
% Fig.7 -- Qualitative (visual) comparison at x4 on hard samples.
% outline.md §3.4 / data-collection-checklist D3 "Fig.7".
% Two-column figure (figure*); \input into the paper after §3.4.
%
% HARD SAMPLES (selected 2026-06-27 by GT high-frequency energy ranking +
%   SR-literature convention; see paper/scripts/build_fig7_assets.py):
%     Urban100 img_092, Urban100 img_024  (dense building grids)
%     Manga109 ThatsIzumiko_000           (high-contrast line art)
%
% ASSET PROVENANCE (paper/figures/fig7_assets/, generated locally):
%   - GT, LR  : CATANet/datasets/TestDataSR/{HR,LR/LRBI}/.../x4  (local datasets)
%   - Bicubic : LR upscaled x4 (PIL BICUBIC) -- real, computed locally
%   - DPRNet  : our SR output ("_CATANet" is only the arch registry name).
%   - Competitors (IMDN, SwinIR-light, SRFormer-light, CATANet): real SR outputs
%     produced on the training machine, cropped to the SAME fixed box by
%     paper/scripts/build_fig7_competitor_crops.py. (RFDN dropped: its AIM x4
%     weights gave sub-bicubic results under our LRBI protocol.)
%
% NOTE: per-image PSNR/SSIM intentionally NOT printed on the crops -- this is a
%   purely qualitative comparison; quantitative numbers live in the main table.
% =====================================================================
\newcommand{\figsevenrow}[2]{%
  % #1 = asset tag   #2 = row caption (shown under the GT full image)
  \begin{subfigure}{0.12\textwidth}
    \includegraphics[width=\linewidth]{figures/fig7_assets/#1_GT_full_box.png}
    \caption*{\tiny #2}
  \end{subfigure}\hfill
  \begin{subfigure}{0.12\textwidth}
    \includegraphics[width=\linewidth]{figures/fig7_assets/#1_Bicubic_crop.png}
    \caption*{\tiny Bicubic}
  \end{subfigure}\hfill
  \begin{subfigure}{0.12\textwidth}
    \includegraphics[width=\linewidth]{figures/fig7_assets/#1_IMDN_crop.png}
    \caption*{\tiny IMDN}
  \end{subfigure}\hfill
  \begin{subfigure}{0.12\textwidth}
    \includegraphics[width=\linewidth]{figures/fig7_assets/#1_SwinIR-light_crop.png}
    \caption*{\tiny SwinIR-light}
  \end{subfigure}\hfill
  \begin{subfigure}{0.12\textwidth}
    \includegraphics[width=\linewidth]{figures/fig7_assets/#1_SRFormer-light_crop.png}
    \caption*{\tiny SRFormer-light}
  \end{subfigure}\hfill
  \begin{subfigure}{0.12\textwidth}
    \includegraphics[width=\linewidth]{figures/fig7_assets/#1_CATANet_crop.png}
    \caption*{\tiny CATANet}
  \end{subfigure}\hfill
  \begin{subfigure}{0.12\textwidth}
    \includegraphics[width=\linewidth]{figures/fig7_assets/#1_DPRNet_crop.png}
    \caption*{\tiny DPRNet (ours)}
  \end{subfigure}\hfill
  \begin{subfigure}{0.12\textwidth}
    \includegraphics[width=\linewidth]{figures/fig7_assets/#1_GT_crop.png}
    \caption*{\tiny GT}
  \end{subfigure}%
}

\begin{figure*}[t]
\centering
\figsevenrow{Urban100_img_092}{Urban100: img\_092}

\vspace{2pt}
\figsevenrow{Urban100_img_024}{Urban100: img\_024}

\vspace{2pt}
\figsevenrow{Manga109_ThatsIzumiko_000}{Manga109: ThatsIzumiko}

\caption{Qualitative comparison at $\times4$ on texture-dense hard samples. The
leftmost panel shows the ground-truth image with the (red) crop region; the
remaining panels magnify that region for Bicubic, IMDN, SwinIR-light,
SRFormer-light, CATANet, DPRNet (ours), and the ground truth. DPRNet restores
repetitive building grids and high-contrast line art with sharper, less aliased
structure, on par with the strongest competitors. This is a qualitative
comparison; quantitative PSNR/SSIM are reported in the main table.}
\label{fig:visual}
\end{figure*}
